%==============================================================================
% Sjabloon poster bachproef
%==============================================================================
% Gebaseerd op document class `a0poster' door Gerlinde Kettl en Matthias Weiser
% Aangepast voor gebruik aan HOGENT door Jens Buysse en Bert Van Vreckem

\documentclass[a0,portrait]{hogent-poster}
\usepackage{enumitem}

% Info over de opleiding
\course{Bachelorproef}
\studyprogramme{toegepaste informatica}
\academicyear{2025-2026}
\institution{Hogeschool Gent, Valentin Vaerwyckweg 1, 9000 Gent}

% Info over de bachelorproef
\title{Een analyse van REXX en Python voor automatisatie op z/OS.}
%\subtitle{Ondertitel (eventueel)}
\author{Vincent Gielen}
\email{vincent.gielen@student.hogent.be}
\supervisor{Dhr. L. Blondeel}
\cosupervisor{Dhr. B. Tjassens Keiser (ICU/Rabobank)}

% Indien ingevuld, wordt deze informatie toegevoegd aan het einde van de
% abstract. Zet in commentaar als je dit niet wilt.
\specialisation{Mainframe}
\keywords{Mainframe, Modernisatie, Automatisatie, Python Rexx}
\projectrepo{https://github.com/Vincent-Gielen/Bachelorproef-Vincent-Gielen}

\begin{document}

\maketitle

\begin{abstract}
Deze bachelorproef gaat over automatisatie op het mainframe, meer specifiek een vergelijking en analyse van Rexx en Python. Rexx bestaat al heel lang op het mainframe, terwijl Python in vergelijking een nieuwkomer is. Door de skills gap en door de ontwikkeling van nieuwe technologieën buiten het mainframe, wordt er echter geargumenteerd dat Python Rexx kan vervangen als meest gebruikte automatisatie taal op het mainframe. Het doel van dit onderzoek is achterhalen hoe Rexx en Python zich qua automatisatie op z/OS verhouden vanuit technisch en praktisch oogpunt.  
Uit het onderzoek blijkt dat Rexx en Python elk hun eigen kwaliteiten hebben. Python wordt steeds populairder, en wordt meer en meer ondersteund op het mainframe, waardoor verschillende use cases van Rexx kunnen worden overgenomen. Rexx is echter zo verweven met de tools op het mainframe dat het (voorlopig) nog steeds onmisbaar is. 
\end{abstract}

\begin{multicols}{2} % This is how many columns your poster will be broken into, a portrait poster is generally split into 2 columns

\section{Verhouding Rexx en Python op z/OS?}

\bigskip

\begin{itemize}
    \item Waarom modernisatie?
    \item Wat doet automatisatie op z/OS?
    \item Voor- en nadelen van Rexx op z/OS?
    \item Voor- en nadelen van Python op z/OS?
    \item Rexx vervangbaar door Python?
\end{itemize}

\bigskip

\textbf {Resultaten literatuurstudie:}
\begin{center}
    \captionsetup{type=figure}
    \includegraphics[width=0.6\linewidth]{req1.png}
    \captionof{figure}{Tabel met de vergelijking van Rexx en Python uit de literatuurstudie}
\end{center}


\section{Methode}

\bigskip

Uit de interviews en de enquête zijn use cases en criteria gehaald.

\bigskip

\textbf{Ervaringsniveau respondenten van de enquête:}
\begin{center}
    \captionsetup{type=figure}
    \includegraphics[width=1.0\linewidth]{exp2.png}
    \captionof{figure}{Ervaringsniveau met z/OS, REXX en Python respondenten}
\end{center}

\bigskip

\textbf{Use cases:}
\begin{itemize}[leftmargin=1.5in]
    \item PoC 1: Gebruik maken van een API
    \item PoC 2: Genereren van JCL
    \item PoC 3: Gebruik maken van RACF
\end{itemize}

\bigskip

\textbf{Requirements:}
\begin{center}
    \captionsetup{type=figure}
    \includegraphics[width=0.6\linewidth]{req3.png}
    \captionof{figure}{De requirements gehaald uit de enquête en interviews}
\end{center}

\section{Resultaten}

\bigskip

\subparagraph{Resultaten PoC's:}
\begin{center}
    \captionsetup{type=figure}
    \includegraphics[width=0.7\linewidth]{req2.png}
    \captionof{figure}{Tabel met de vergelijking van Rexx en Python uit de resultaten van de PoC's}
\end{center}

\section{Conclusie}
\bigskip
\begin{itemize}
    \setlength\itemsep{0.5em}
    \item Rexx en Python hebben \textbf{beiden een plaats} op het mainframe
    \item \textbf{Rexx} brengt stabiliteit, veiligheid en integratie met tools op het mainframe
    \item \textbf{Python} is nodig voor integratie met technologieën buiten het mainframe en heeft veel nieuwe functionaliteiten en ondersteuning
    \item Voor bepaalde \textbf{specifieke use cases} kan de ene taal een betere oplossing bieden dan de andere
    \item Voor andere use cases is de keuze tussen beide talen veelal gebaseerd op \textbf{voorkeur en expertise}
    \item Aanbeveling modernisatie strategie voor bedrijven:
\end{itemize}

\bigskip

\begin{center}
    \captionsetup{type=figure}
    \includegraphics[width=0.6\linewidth]{cirkelconclusie.png}
    \captionof{figure}{Visualisatie aanbeveling modernisatie strategie}
\end{center}

\bigskip

\section{Toekomstig onderzoek}
\bigskip
\begin{itemize}
    \setlength\itemsep{0.5em}
    \item \textbf{Beperkingen:} Beter systeem en meer tijd zijn nodig
    \item \textbf{Breder onderzoek:} Meer technologieën, use cases en requirements
    \item \textbf{Toekomst} van automatisatie
\end{itemize}


\end{multicols}
\end{document}